\lyseis{Β3.9}{Άθροισμα γωνιών τριγώνου}
\sol{1} α) Λάθος, οι δύο ορθές κάνουν ήδη $180^\circ$ και η τρίτη γωνία θα ήταν $0^\circ$. β) Σωστό, $50^\circ+60^\circ+70^\circ=180^\circ$. γ) Σωστό, αφού $\hat{\mathrm{Β}}+\hat{\mathrm{Γ}}=180^\circ-90^\circ=90^\circ$. δ) Λάθος, κάθε γωνία του ισόπλευρου είναι $60^\circ$. ε) Σωστό, $(180^\circ-30^\circ):2=75^\circ$. στ) Λάθος, δύο αμβλείες θα είχαν άθροισμα πάνω από $180^\circ$.
\sol{2} α) $180^\circ-(45^\circ+65^\circ)=70^\circ$. β) $180^\circ-(90^\circ+34^\circ)=56^\circ$. γ) $180^\circ-(108^\circ+25^\circ)=47^\circ$. δ) $180^\circ-(59^\circ+61^\circ)=60^\circ$.
\sol{3} Κάθε γωνία βάσης είναι $(180^\circ-\text{κορυφή}):2$. α) $70^\circ$. β) $40^\circ$. γ) $45^\circ$. δ) $80^\circ$.
\sol{4} Η κορυφή είναι $180^\circ-2\cdot(\text{βάση})$. α) $50^\circ$. β) $120^\circ$. γ) $90^\circ$. δ) $20^\circ$.
\sol{5} $x+(x+10^\circ)+(x+20^\circ)=180^\circ$, άρα $3x+30^\circ=180^\circ$, $3x=150^\circ$ και $x=50^\circ$. Οι γωνίες είναι $50^\circ$, $60^\circ$ και $70^\circ$.
\sol{6} Μπορούν οι τριάδες με άθροισμα $180^\circ$: α) $180^\circ$, ναι. β) $190^\circ$, όχι. γ) $180^\circ$, ναι. δ) $190^\circ$, όχι. ε) $180^\circ$, ναι. στ) $170^\circ$, όχι.
\sol{7} α) οξυγώνιο (όλες οξείες). β) ορθογώνιο (έχει μια γωνία $90^\circ$). γ) αμβλυγώνιο (έχει γωνία $105^\circ$). δ) οξυγώνιο.
\sol{8} Κάθε γωνία της βάσης είναι $(180^\circ-110^\circ):2=35^\circ$.
\sol{9} Αν η $50^\circ$ είναι η κορυφή: οι γωνίες της βάσης είναι $(180^\circ-50^\circ):2=65^\circ$, δηλαδή $(50^\circ,65^\circ,65^\circ)$. Αν η $50^\circ$ είναι γωνία της βάσης: η άλλη της βάσης είναι $50^\circ$ και η κορυφή $180^\circ-100^\circ=80^\circ$, δηλαδή $(80^\circ,50^\circ,50^\circ)$.
\sol{10} Ο τοίχος και το δάπεδο σχηματίζουν ορθή γωνία, άρα το τρίγωνο είναι ορθογώνιο και οι οξείες γωνίες του έχουν άθροισμα $90^\circ$. Η γωνία με τον τοίχο είναι $90^\circ-72^\circ=18^\circ$.
\sol{11} $2x+3x+4x=9x=180^\circ$, άρα $x=20^\circ$. Οι γωνίες είναι $40^\circ$, $60^\circ$, $80^\circ$ και το τρίγωνο είναι οξυγώνιο.
\sol{12} Αν η γωνία της βάσης είναι $\beta$, η κορυφή είναι $3\beta$ και $\beta+\beta+3\beta=180^\circ$, άρα $5\beta=180^\circ$ και $\beta=36^\circ$. Οι γωνίες είναι $36^\circ$, $36^\circ$ και $108^\circ$.
\sol{13} Αν είχε δύο γωνίες όχι οξείες, καθεμία τουλάχιστον $90^\circ$, το άθροισμά τους θα ήταν τουλάχιστον $180^\circ$ και η τρίτη γωνία θα ήταν $0^\circ$ ή αρνητική, που δεν γίνεται. Άρα το πολύ μία γωνία δεν είναι οξεία και οι άλλες δύο είναι οξείες.
\sol{14} Οι γωνίες της βάσης είναι $\hat{\mathrm{Β}}=\hat{\mathrm{Γ}}=(180^\circ-44^\circ):2=68^\circ$. Επειδή $\mathrm{ΔΕ}\parallel\mathrm{ΒΓ}$, είναι $\widehat{\mathrm{ΑΔΕ}}=\hat{\mathrm{Β}}=68^\circ$ και $\widehat{\mathrm{ΑΕΔ}}=\hat{\mathrm{Γ}}=68^\circ$ (εντός--εκτός και επί τα αυτά). Η $\widehat{\mathrm{ΔΕΓ}}$ είναι παραπληρωματική της $\widehat{\mathrm{ΑΕΔ}}$, άρα $180^\circ-68^\circ=112^\circ$.
\sol{15} Οι διχοτόμοι χωρίζουν τις γωνίες σε $35^\circ$ και $55^\circ$. Στο τρίγωνο $\mathrm{ΑΚΒ}$ είναι $\widehat{\mathrm{ΑΚΒ}}=180^\circ-35^\circ-55^\circ=90^\circ$. Ναι, βγαίνει το ίδιο πάντα: αν η μία γωνία είναι $\omega$, η άλλη είναι $180^\circ-\omega$, και οι μισές τους έχουν άθροισμα $\frac{\omega}{2}+\frac{180^\circ-\omega}{2}=90^\circ$, άρα η γωνία στο $\mathrm{Κ}$ είναι πάντα $90^\circ$.
\sol{16} Αν $\hat{\Gamma}=y$ τότε $\hat{\mathrm{Α}}=4y$ και $4y+30^\circ+y=180^\circ$, άρα $5y=150^\circ$ και $y=30^\circ$. Οι γωνίες είναι $\hat{\mathrm{Α}}=120^\circ$, $\hat{\mathrm{Β}}=30^\circ$, $\hat{\mathrm{Γ}}=30^\circ$. Είναι αμβλυγώνιο ως προς τις γωνίες και ισοσκελές ως προς τις πλευρές (έχει δύο ίσες γωνίες, με κορυφή το $\mathrm{Α}$).
\sol{17} Ελεύθερη εργασία. Οι τρεις γωνίες, όταν μπουν δίπλα δίπλα, σχηματίζουν ευθεία γωνία ($180^\circ$), και στο αμβλυγώνιο τρίγωνο επίσης.
\sol{18} Ελεύθερη κατασκευή. $\hat{\mathrm{Α}}=180^\circ-50^\circ-60^\circ=70^\circ$ και η μέτρηση θα πρέπει να δώσει περίπου $70^\circ$.
